\section{Complexity and the Security--Cost Frontier}
\label{sec:complexity}

\Aegis buys robustness by running a committee, and the idea is explicit that the cost of
that committee may be unattractive at small $n$. We make the cost quantitative and pair it
with the tolerance it purchases, so the trade-off is stated rather than hidden.

\subsection{Token and API cost}
Each request issues one hardened-judge call per committee member plus the Analyzer calls.
Let $\ell_{\text{in}}$ and $\ell_{\text{out}}$ be the per-judge input and output token
counts. The token cost per adjudication is
\begin{equation}
  \text{Tokens}(n) \;=\; n\,(\ell_{\text{in}}+\ell_{\text{out}}) \;+\; \text{Tokens}_{\text{analyze}},
  \label{eq:tokencost}
\end{equation}
i.e.\ \emph{linear} in the committee size $n$, since judges are independent and the payload
is broadcast unchanged. Payload isolation adds only the fixed delimiter overhead. This
linearity is the cost the idea asks to quantify, and it is the reason $n$ is restricted to
the range $1$--$7$: the token bill grows in lockstep with $n$ while, by \cref{lem:gmedrobust},
the tolerated fault count grows only as $\lfloor (n-1)/2\rfloor$ (median rules) or
$\lfloor (n-2)/2\rfloor$ (Krum).

\subsection{Latency}
In the parallel topology the judge calls overlap, so wall-clock latency is
\begin{equation}
  \text{Lat}_{\parallel}(n) \;=\; \max_{k\in[n]} \text{lat}(J_k) \;+\; \text{lat}(\Agg) \;+\; \text{lat}_{\text{analyze}},
  \label{eq:latpar}
\end{equation}
dominated by the slowest judge plus a small aggregation term; in the sequential topology it
is $\sum_k \text{lat}(J_k)+\text{lat}(\Agg)$, linear in $n$. The idea notes latency is
``linear in the number of agents''; \cref{eq:latpar} shows this is the \emph{sequential}
worst case and that parallel deployment reduces the judge contribution to a constant in $n$
(at unchanged token cost), trading hardware concurrency for latency.

\subsection{Aggregation cost}
The three rules cost, per request, $O(nd)$ with an $O(n\log n)$ per-coordinate selection for
$\cmed$; $O(Tnd)$ for $\gmed$ with $T$ Weiszfeld iterations~\cite{pillutla2022rfa}; and
$O(n^2 d)$ for $\Krum$ from its pairwise distances~\cite{blanchard2017krum}. With $n\le 7$
and modest embedding dimension $d=1+m$, all three are dominated by the LLM calls of
\cref{eq:tokencost} and contribute negligibly to latency.

\subsection{The frontier}
Define the \emph{Byzantine tolerance fraction} $\alpha^\star(n)=\lfloor(n-1)/2\rfloor/n$ for
median rules. \Cref{tab:cost} tabulates tolerance against the multiplicative token cost
$n$. The security gain is discrete and, at small $n$, coarse: moving from $n=3$ (tolerates
$f=1$) to $n=5$ (tolerates $f=2$) doubles neither cost-effectiveness nor tolerance cleanly.
Whether a given operating point is worthwhile is an empirical question about the deployment's
threat level and latency budget, which \cref{sec:experiments} is designed to answer rather
than to presuppose.

\begin{table}[t]
\centering
\caption{Security--cost frontier for median-type aggregation. Tolerance is the maximum $f$
with $f<n/2$; token cost is the multiple of a single-judge call from \cref{eq:tokencost}
(Analyzer overhead excluded). Small committees tolerate very few faults, the feasibility
risk the idea flags as an open empirical question.}
\label{tab:cost}
\small
\begin{tabular}{@{}cccc@{}}
\toprule
Committee $n$ & Max tolerated $f$ & Tolerance fraction $\alpha^\star$ & Token cost (\(\times\) single judge) \\
\midrule
1 & 0 & $0.00$ & $1$ \\
2 & 0 & $0.00$ & $2$ \\
3 & 1 & $0.33$ & $3$ \\
4 & 1 & $0.25$ & $4$ \\
5 & 2 & $0.40$ & $5$ \\
6 & 2 & $0.33$ & $6$ \\
7 & 3 & $0.43$ & $7$ \\
\bottomrule
\end{tabular}
\end{table}
